% !TEX TS-program = pdflatexmk
% !TEX program = pdflatex
% Derived from problems/template.tex; build with make problem HW=hw01.
\documentclass[nohyper,nobib,nofonts]{tufte-handout}
\newcommand{\bookroot}{..}
\IfFileExists{preamble.tex}{\renewcommand{\bookroot}{.}}{}
\makeatletter
\edef\input@path{{\bookroot/}}
\makeatother
% Shared packages and typography for the book and homework handouts.
% Handouts set this directory when compiled from problems/.
\providecommand{\bookroot}{.}
\usepackage[T1]{fontenc}
% Give titletoc a separate part number before nameref wraps \@part.
\usepackage{etoolbox}
\makeatletter
\@ifclassloaded{tufte-book}{%
  \patchcmd{\@part}
    {\thepart\hspace{1em}}
    {\protect\numberline{\thepart}}
    {}
    {\PackageError{linear-algebra-notes}{Could not format part contents labels}
      {Check the book class definition of \string\@part.}}
}{}
\makeatother

\usepackage{nameref}
\usepackage{amsmath,amssymb,amsthm}
\usepackage{aliascnt} % Shared theorem numbering with distinct autoref names.
\usepackage{graphicx}
\usepackage{cancel}
\usepackage{tikz}
\usetikzlibrary{arrows.meta,calc}
\usepackage{booktabs}
\usepackage{multicol} % Two-column appendix formula sheet.
\usepackage{needspace} % Keep an exercise heading with its opening lines.
\usepackage{makeidx}
\usepackage{url}
% The same alphabetic AMS style as Maths Analysis, bundled for Overleaf.
% Uncited homework exports need neither the book database nor a Biber pass.
\makeatletter
\@ifclassloaded{tufte-book}{%
  \ifdefined\standaloneformulasheet\else
    \providecommand{\input@path}{}
    \edef\input@path{{\bookroot/styles/}%
      {\bookroot/styles/biblatex-ams/bbx/}%
      {\bookroot/styles/biblatex-ams/cbx/}\input@path}
    \usepackage[backend=biber,natbib=true,bibstyle=analysis-ams,
      citestyle=ams-alphabetic,sorting=nyt]{biblatex}
    \renewcommand*{\multicitedelim}{\addcomma\space}
    \expandafter\addbibresource\expandafter{\bookroot/bibliography.bib}
  \fi
}{}
\makeatother

\graphicspath{{\bookroot/figures/}}
\setkeys{Gin}{width=\linewidth,totalheight=\textheight,keepaspectratio}
\setcounter{secnumdepth}{1}

% Retain the borrowed title/contents design only for the book.
% Adapted from Kevin Godby's Tufte-LaTeX title/contents example:
% https://groups.google.com/forum/#!topic/tufte-latex/ujdzrktC1BQ
\makeatletter
\@ifclassloaded{tufte-book}{%
  \makeindex
  % Leave room for two-digit appendix figure numbers such as C.13.
  \patchcmd{\@tufte@lof@line}{\@tempdima 2.0em}{\@tempdima 3.0em}
    {}
    {\PackageError{linear-algebra-notes}{Could not widen list-of-figures numbers}
      {Check the Tufte list-of-figures definition.}}
  \setcounter{tocdepth}{0} % Parts and chapters only.
  % Title-only chapter headings; retain counters for sections and references.
  \titleformat{\chapter}[block]
    {\relax\ifthenelse{\NOT\boolean{@tufte@symmetric}}{\begin{fullwidth}}{}}
    {}
    {0pt}
    {\huge\rmfamily\itshape}
    [\ifthenelse{\NOT\boolean{@tufte@symmetric}}{\end{fullwidth}}{}%
      % Clear the fullwidth list's pending paragraph cleanup so it cannot
      % overwrite the chapter hook and suppress the next section's spacing.
      \@endpefalse]
  \renewcommand{\maketitlepage}{%
    \begingroup
    \setlength{\parindent}{0pt}
    {\fontsize{24}{24}\selectfont\textit{\@author}\par}
    \vspace{1.75in}
    \begin{fullwidth}
      \raggedright
      {\fontsize{36}{44}\selectfont\@title\par}
    \end{fullwidth}
    \vspace{0.5in}
    {\fontsize{14}{14}\selectfont\textsf{\smallcaps{\@date}}\par}
    \vfill
    {\fontsize{14}{14}\selectfont\textit{\@publisher}\par}
    \thispagestyle{empty}
    \endgroup
  }
  \titlecontents{part}
    [0pt]
    {\addvspace{0.75\baselineskip}\normalfont\normalsize}
    {\allcaps{Part~\thecontentslabel}\quad\allcaps}
    {\allcaps}
    {}
    [\vspace{0.5\baselineskip}]
  % Hide chapter numbers in the contents as well.
  \titlecontents{chapter}
    [4em]
    {\normalfont\normalsize}
    {\textit}
    {\textit}
    {\qquad\thecontentspage}
    [\vspace{0.5\baselineskip}]
}{}
\makeatother

% Enable links after the Tufte contents customizations in both the book and
% handouts. Their nohyper option defers loading; keep sidenotes unlinked.
\usepackage[hidelinks,hyperfootnotes=false,bookmarksdepth=0]{hyperref}
\makeatletter
\pdfstringdefDisableCommands{\let\thanklessauthor\@author}
\makeatother

% Match the ordinary Tufte margin text. Only the label is bold; the full
% stop separates it from regular serif text, with no first-line indent.
% Tufte continues to place figures and their captions in its native layout.
\usepackage{caption}
\captionsetup{
  format=plain,
  font={normalfont,footnotesize},
  labelfont=bf,
  textfont=normalfont,
  labelsep=period,
  justification=RaggedRight,
  singlelinecheck=false,
  parindent=0pt,
  skip=0pt
}
% Tufte assembles all captions in separate boxes; link to the caption.
\captionsetup{hypcap=false}

% Schedule complete margin notes and figures, as in Maths Analysis.
\usepackage{marginfix}
\setlength{\marginparpush}{12pt}
\setlength{\marginheightadjustment}{-6pt}
\makeatletter
\@ifclassloaded{tufte-book}{%
  \AtBeginDocument{%
    \BeforeBeginEnvironment{exercise}{\Needspace{3\baselineskip}}%
  }
}{}
% Offsets are preferred positions rather than hidden space in measured boxes.
\renewcommand{\marginnote}[2][0pt]{%
  \ifthenelse{\boolean{@tufte@loadnatbib}}{%
    \let\cite\@tufte@infootnote@cite
  }{}%
  \gdef\@tufte@citations{}%
  \mparshift{#1}%
  \marginpar{%
    \@tufte@marginnote@font\@tufte@marginnote@justification
    \@tufte@margin@par\noindent #2%
  }%
  \@tufte@print@citations
  \ifthenelse{\boolean{@tufte@loadnatbib}}{%
    \let\cite\@tufte@normal@cite
  }{}%
}
\renewenvironment{@tufte@margin@float}[2][-1.2ex]
  {\FloatBarrier
    \def\book@margin@offset{#1}%
    \begin{lrbox}{\@tufte@margin@floatbox}%
      \begin{minipage}[t]{\marginparwidth}%
        \@tufte@caption@font
        \def\@captype{#2}%
        \@tufte@caption@justification
        \@tufte@margin@par\noindent
  }
  {\end{minipage}%
    \end{lrbox}%
    \ifdim\dimexpr\ht\@tufte@margin@floatbox+\dp\@tufte@margin@floatbox\relax
        <\dimexpr\textheight+\marginheightadjustment\relax
      \mparshift{\book@margin@offset}%
      \marginpar{\usebox{\@tufte@margin@floatbox}}%
    \else
      \PackageError{linear-algebra-notes}{Margin float is taller than the usable margin}
        {Use a main-text figure or table instead; keep labels readable.}%
    \fi
  }
% Ordinary and full-width floats retain Tufte's standard height handling.
\makeatother

% Shared notation and environments. Keep topic-specific macros local.
\newcommand{\N}{\mathbb{N}}
\newcommand{\Z}{\mathbb{Z}}
\newcommand{\Q}{\mathbb{Q}}
\newcommand{\R}{\mathbb{R}}
\newcommand{\C}{\mathbb{C}}

% Reusable draft styling. Existing author-requested red reviews stay red.
\definecolor{pendingreviewcolor}{HTML}{486581}
\newenvironment{pendingreview}{\par\color{pendingreviewcolor}}{\par}
\newcommand{\pendingreviewtext}[1]{\textcolor{pendingreviewcolor}{#1}}

% The same linked reference typography as Maths Analysis.
\renewcommand*{\figureautorefname}{Figure}
\renewcommand*{\tableautorefname}{Table}
\renewcommand*{\partautorefname}{Part}
\renewcommand*{\chapterautorefname}{Chapter}
\renewcommand*{\appendixautorefname}{Appendix}
\renewcommand*{\sectionautorefname}{\S}
\renewcommand*{\subsectionautorefname}{\S}
\renewcommand*{\subsubsectionautorefname}{\S}
\renewcommand*{\theoremautorefname}{Theorem}
\def\equationautorefname~#1\null{Equation~(#1)\null}

% Keep submission-specific layout hints out of shared book exercises.
\makeatletter
\newcommand{\handoutonly}[1]{\@ifclassloaded{tufte-handout}{#1}{}}
% Append after the complete statement and hints, before the exercise ends.
% The handouts print solutions immediately; only the book needs these links.
\newcommand{\solutionlink}[1]{%
  \@ifclassloaded{tufte-book}{%
    \ifhmode\unskip\space\else\noindent\fi
    \textit{\hyperref[sol:#1]{Solution on \autopageref*{sol:#1}.}}%
  }{}%
}
\makeatother

% Select the original exercise number for the next appendix restatement.
\newcommand{\restatedexercisekey}{}
\newcommand{\solutionheading}[1]{%
  \renewcommand{\restatedexercisekey}{#1}%
}

% Reuse the exercise's theorem styling and original number. Register only
% the solution anchor, not a second exercise label or solution link.
\newenvironment{restatedproblem}{%
  \let\restatedproblemlabel\label
  \renewenvironment{exercise}[1][]{%
    \ifstrempty{##1}{%
      \begin{restatedexercise}%
    }{%
      \begin{restatedexercise}[##1]%
    }%
    \phantomsection
    \restatedproblemlabel{sol:\restatedexercisekey}%
  }{\end{restatedexercise}}%
  \renewcommand{\label}[1]{}%
  \renewcommand{\solutionlink}[1]{}%
}{\par}

% Publication month and year used on the copyright page.
\newcommand{\monthyear}{%
  \ifcase\month\or January\or February\or March\or April\or May\or June\or
  July\or August\or September\or October\or November\or December\fi
  \space\number\year
}

% One shared counter makes results unambiguous within each section.
% Alias counters let autoref identify each environment without renumbering it.
\theoremstyle{plain}
\newtheorem{theorem}{Theorem}[section]
\newaliascnt{lemma}{theorem}
\newtheorem{lemma}[lemma]{Lemma}
\aliascntresetthe{lemma}
\newaliascnt{proposition}{theorem}
\newtheorem{proposition}[proposition]{Proposition}
\aliascntresetthe{proposition}
\newaliascnt{corollary}{theorem}
\newtheorem{corollary}[corollary]{Corollary}
\aliascntresetthe{corollary}
\theoremstyle{definition}
\newaliascnt{definition}{theorem}
\newtheorem{definition}[definition]{Definition}
\aliascntresetthe{definition}
\newaliascnt{example}{theorem}
\newtheorem{example}[example]{Example}
\aliascntresetthe{example}
\newaliascnt{exercise}{theorem}
\newtheorem{exercise}[exercise]{Exercise}
\aliascntresetthe{exercise}
\newtheorem*{restatedexercise}{\autoref{ex:\restatedexercisekey}}
\theoremstyle{remark}
\newaliascnt{remark}{theorem}
\newtheorem{remark}[remark]{Remark}
\aliascntresetthe{remark}
\providecommand*{\lemmaautorefname}{Lemma}
\providecommand*{\propositionautorefname}{Proposition}
\providecommand*{\corollaryautorefname}{Corollary}
\providecommand*{\definitionautorefname}{Definition}
\providecommand*{\exampleautorefname}{Example}
\providecommand*{\exerciseautorefname}{Exercise}
\providecommand*{\remarkautorefname}{Remark}
\providecommand*{\homeworkproblemautorefname}{Problem}

% Build the theorem catalogue from existing numbered theorem headings.
\makeatletter
\@ifclassloaded{tufte-book}{%
  \newcommand{\listoftheorems}{%
    \begingroup
    \let\cleardoublepage\clearpage
    \chapter*{List of Theorems}%
    \@starttoc{loe}%
    \endgroup
  }
  \let\l@theorem\@tufte@lof@line
  \providecommand*{\toclevel@theorem}{2}
}{}
\ifdefined\listoftheorems
  \pretocmd{\@begintheorem}{%
    \ifstrequal{#1}{Theorem}{%
      \ifblank{#2}{}{%
        \addcontentsline{loe}{theorem}{\protect\numberline{#2}#3}%
      }%
    }{}%
  }{}{%
    \PackageError{linear-algebra-notes}{Could not register theorem list entries}
      {Check the amsthm definition of \string\@begintheorem.}%
  }
\fi
\makeatother


\counterwithout{theorem}{section}
\renewcommand{\thetheorem}{\arabic{theorem}}
\theoremstyle{definition}
\newaliascnt{homeworkproblem}{theorem}
\newtheorem{homeworkproblem}[homeworkproblem]{Problem}
\aliascntresetthe{homeworkproblem}
\renewenvironment{exercise}{\begin{homeworkproblem}}{\end{homeworkproblem}}

\title[MA-GY7033 Linear Algebra I, Homework 1]{MA-GY7033 Linear Algebra I\\Homework 1}
\author[Lucas Barbosa, lc6424]{Lucas Barbosa, \texttt{lc6424}}
\date{}
\fancyhead{}
\fancyhead[LE,RO]{\footnotesize MA-GY7033 Linear Algebra I, Homework 1\quad\thepage}

\begin{document}
\maketitle

% Source: supplied Homework 1 PDF, p. 1, assignment-wide conventions.
Throughout, \(\mathbb F\) denotes \(\R\) or \(\C\), and vector spaces
are over \(\mathbb F\) unless another field is specified. The symbol
\(\subseteq\) allows equality. The space \(\mathcal P_n(\mathbb F)\)
consists of all polynomials of degree at most \(n\) with coefficients in
\(\mathbb F\), including the zero polynomial. Operations on functions
and sequences are pointwise unless stated otherwise.

% Replace the source's lecture-PDF filename with self-contained topic wording.
Justify every answer. Use the definitions and results on vector spaces and
bases, together with elementary algebra and any results you have proved in
earlier problems. Matrices, determinants, and theorems about linear maps
are not needed. The final two problems are challenges.

\section*{Axioms, subspaces, and span}
% !TEX root = ../../hw01.tex
% Source: ma-gy7033fall2026homework1.pdf, p. 1, Problem 1(a--d).
% The assignment-wide choice F=R or C is made explicit for shared reuse.
\begin{exercise}[Consequences of the axioms]
  \label{ex:hw01:axiom-consequences}
  Let \(\mathbb F\) denote \(\R\) or \(\C\). Let \(V\) be a vector
  space over \(\mathbb F\), \(v\in V\), and \(a\in\mathbb F\).
  Prove directly from the vector space axioms that:
  \begin{enumerate}[(a)][2]
    \item the zero vector and the additive inverse of each vector are unique;
    \item \(0v=0_V\) and \(a0_V=0_V\);
    \item \((-1)v=-v\);
    \item if \(av=0_V\), then \(a=0\) or \(v=0_V\).
  \end{enumerate}
  Distinguish scalar zeros from vector zeros in your proofs.
  \solutionlink{hw01:axiom-consequences}
\end{exercise}

% !TEX root = ../../hw01.tex
\noindent\emph{Awaiting solution.}


% !TEX root = ../../hw01.tex
% Source: ma-gy7033fall2026homework1.pdf, p. 1, Problem 2(a--b).
\begin{exercise}[An unfamiliar vector space]
  \label{ex:hw01:positive-real-vector-space}
  On \(V=\{x\in\R:x>0\}\), define
  \[
    x\oplus y=xy,\qquad a\odot x=x^a\qquad(a\in\R).
  \]
  \begin{enumerate}[(a)][2]
    \item Prove that these operations make \(V\) a vector space over
      \(\R\). Identify its zero vector and the additive inverse of \(x\).
    \item Determine the dimension of \(V\) and characterize all its bases.
  \end{enumerate}
  You may use the usual properties of real powers and logarithms.
  \solutionlink{hw01:positive-real-vector-space}
\end{exercise}

% !TEX root = ../../hw01.tex
\noindent\emph{Awaiting solution.}


% !TEX root = ../../hw01.tex
% Source: ma-gy7033fall2026homework1.pdf, p. 1, Problem 3(a--b).
% Retain the source's unused T and arbitrary, possibly infinite, subsets.
\begin{exercise}[What span does]
  \label{ex:hw01:span-and-containing-subspaces}
  Let \(V\) be a vector space over \(\mathbb F\), where
  \(\mathbb F\) denotes \(\R\) or \(\C\). Let \(S,T\subseteq V\)
  be arbitrary subsets, possibly empty or infinite. Recall that an empty
  linear combination is equal to \(0_V\).
  \begin{enumerate}[(a)][2]
    \item Prove that \(\operatorname{span}(S)\) is a subspace containing \(S\).
    \item Prove that every subspace containing \(S\) contains
      \(\operatorname{span}(S)\). Deduce that \(\operatorname{span}(S)\)
      is the intersection of all subspaces containing \(S\).
  \end{enumerate}
  \solutionlink{hw01:span-and-containing-subspaces}
\end{exercise}

% !TEX root = ../../hw01.tex
\noindent\emph{Awaiting solution.}


% !TEX root = ../../hw01.tex
% Source: ma-gy7033fall2026homework1.pdf, p. 2, Problem 4(a--b).
\begin{exercise}[A parameter that changes the span]
  \label{ex:hw01:parameter-dependent-span}
  Let \(\mathbb F\) denote \(\R\) or \(\C\). For \(a\in\mathbb F\), set
  \[
    v_1=(1,0,1),\qquad v_2=(0,1,1),\qquad v_3=(1,1,a).
  \]
  \begin{enumerate}[(a)][2]
    \item Determine exactly when \(\{v_1,v_2,v_3\}\) is a basis of
      \(\mathbb F^3\). For each remaining value of \(a\), find a basis
      and the dimension of its span, and describe the span by a linear equation.
    \item When the three vectors form a basis, express any
      \((x,y,z)\in\mathbb F^3\) as their linear combination.
  \end{enumerate}
  \solutionlink{hw01:parameter-dependent-span}
\end{exercise}

% !TEX root = ../../hw01.tex
\noindent\emph{Awaiting solution.}


% !TEX root = ../../hw01.tex
% Source: ma-gy7033fall2026homework1.pdf, p. 2, Problem 5(a--b).
\begin{exercise}[Solutions and differences of solutions]
  \label{ex:hw01:affine-solution-set}
  Let \(\mathbb F\) denote \(\R\) or \(\C\). In \(\mathbb F^3\), define
  \[
    \begin{aligned}
      H&=\{(x,y,z):x+y+z=0,\ x-y=0\},\\
      A&=\{(x,y,z):x+y+z=1,\ x-y=0\}.
    \end{aligned}
  \]
  \begin{enumerate}[(a)][2]
    \item Prove that \(H\) is a subspace and find a basis and its dimension.
      Determine whether \(A\) is a subspace.
    \item Find \(p\in A\) for which \(A=\{p+h:h\in H\}\), and prove
      this equality. Prove that the difference of any two elements of
      \(A\) belongs to \(H\).
  \end{enumerate}
  \solutionlink{hw01:affine-solution-set}
\end{exercise}

% !TEX root = ../../hw01.tex
\noindent\emph{Awaiting solution.}


\section*{Polynomials, functions, and sequences}
% !TEX root = ../../hw01.tex
% Source: ma-gy7033fall2026homework1.pdf, p. 2, Problem 6(a--c).
% Include assignment-wide field and polynomial conventions for shared reuse.
\begin{exercise}[A polynomial determined by three values]
  \label{ex:hw01:polynomial-interpolation}
  Let \(\mathbb F\) denote \(\R\) or \(\C\), and work in
  \(\mathcal P_2(\mathbb F)\), the polynomials of degree at most two
  with coefficients in \(\mathbb F\), including the zero polynomial.
  \begin{enumerate}[(a)][2]
    \item Construct \(q_{-1},q_0,q_1\) such that, for
      \(j,k\in\{-1,0,1\}\),
      \[
        q_j(k)=
        \begin{cases}
          1,&j=k,\\
          0,&j\ne k.
        \end{cases}
      \]
    \item Prove that these polynomials form a basis and that every
      \(p\in\mathcal P_2(\mathbb F)\) satisfies
      \[
        p=p(-1)q_{-1}+p(0)q_0+p(1)q_1.
      \]
    \item Find the unique polynomial \(p\in\mathcal P_2(\mathbb F)\)
      for which \(p(-1)=2\), \(p(0)=-1\), and \(p(1)=4\).
  \end{enumerate}
  \solutionlink{hw01:polynomial-interpolation}
\end{exercise}

% !TEX root = ../../hw01.tex
\noindent\emph{Awaiting solution.}


% !TEX root = ../../hw01.tex
% Source: ma-gy7033fall2026homework1.pdf, p. 3, Problem 7(a--b), including hint.
\begin{exercise}[A vector space of recursively defined sequences]
  \label{ex:hw01:recursive-sequences}
  Let \(\mathbb F\) be \(\R\) or \(\C\). Let \(V\) consist of the
  infinite sequences \((x_0,x_1,\ldots)\), where each \(x_k\in\mathbb F\),
  that satisfy
  \[
    x_{n+2}=x_{n+1}+x_n\qquad(n\ge0).
  \]
  \begin{enumerate}[(a)][2]
    \item Prove that \(V\) is a vector space where vector addition is
      \[
        (x_0,x_1,\ldots)+(y_0,y_1,\ldots)
          =(x_0+y_0,x_1+y_1,\ldots)
      \]
      and scalar multiplication is
      \[
        s(x_0,x_1,\ldots)=(sx_0,sx_1,\ldots).
      \]
      Construct a basis and determine its dimension.
    \item Prove that the only sequence in \(V\) with finitely many nonzero
      terms is the zero sequence.
  \end{enumerate}
  \emph{Hint.} For (a), consider the two sequences with initial values
  \((x_0,x_1)=(1,0)\) and \((x_0,x_1)=(0,1)\).
  \solutionlink{hw01:recursive-sequences}
\end{exercise}

% !TEX root = ../../hw01.tex
\noindent\emph{Awaiting solution.}


\section*{Bases and dimension}
% !TEX root = ../../hw01.tex
% Source: ma-gy7033fall2026homework1.pdf, p. 3, Problem 8(a--b).
\begin{exercise}[Intersecting two spans]
  \label{ex:hw01:intersections-and-sums}
  Let \(\mathbb F\) denote \(\R\) or \(\C\). In \(\mathbb F^4\), let
  \[
    \begin{aligned}
      U&=\operatorname{span}\{(1,0,1,0),(0,1,0,1)\},\\
      W&=\operatorname{span}\{(1,1,0,0),(0,0,1,1)\}.
    \end{aligned}
  \]
  \begin{enumerate}[(a)][2]
    \item Find bases and dimensions of \(U\), \(W\), and \(U\cap W\).
    \item Define \(U+W=\{u+w:u\in U,w\in W\}\). Prove directly that
      \(U+W\) is a subspace, find a basis and its dimension, and describe it
      by one linear equation. Do not assume a dimension formula for sums
      of subspaces.
  \end{enumerate}
  \solutionlink{hw01:intersections-and-sums}
\end{exercise}

% !TEX root = ../../hw01.tex
\noindent\emph{Awaiting solution.}


% !TEX root = ../../hw01.tex
% Source: ma-gy7033fall2026homework1.pdf, p. 3, Problem 9(a--c), including hint.
% Omit only "from the lecture" in the hint for self-contained shared wording.
% Lucas authorized adding the prerequisite comparison theorem on November 9.
\begin{exercise}[Extending an independent set]
  \label{ex:hw01:basis-extension}
  Let \(V\) be a vector space over \(\mathbb F\), where
  \(\mathbb F\) denotes \(\R\) or \(\C\), and let \(V\) have dimension \(n\).
  \begin{enumerate}[(a)][2]
    \item Prove that every linearly independent subset of \(V\) is finite
      and can be extended to a basis of \(V\).
    \item Deduce that any \(n\) linearly independent vectors form a basis,
      and that any set of \(n\) vectors spanning \(V\) is a basis.
    \item Prove that a subspace \(W\subseteq V\) with \(\dim W=n\)
      must equal \(V\). Give an example showing that an infinite dimensional
      space can have a proper subspace of the same dimension; exhibit bases
      to justify the comparison.
  \end{enumerate}
  \emph{Hint.} For (a), show that adjoining \(v\notin\operatorname{span}(S)\)
  to an independent set \(S\) preserves independence. Use the comparison
  theorem to bound the number of steps.
  \solutionlink{hw01:basis-extension}
\end{exercise}

% !TEX root = ../../hw01.tex
\noindent\emph{Awaiting solution.}


% !TEX root = ../../hw01.tex
% Source: ma-gy7033fall2026homework1.pdf, p. 4, Problem 10(a--e).
\begin{exercise}[Testing plausible assertions]
  \label{ex:hw01:span-and-independence-assertions}
  Decide whether each statement is true or false. Prove every true statement
  and give an explicit counterexample for every false one. Here \(S,T\)
  are subsets of a vector space \(V\) over \(\mathbb F\), where
  \(\mathbb F\) denotes \(\R\) or \(\C\).
  \begin{enumerate}[(a)][2]
    \item If \(S\) and \(T\) are linearly independent, then
      \(S\cup T\) is linearly independent.
    \item If \(S\) and \(T\) are linearly independent, then
      \(S\cap T\) is linearly independent.
    \item If \(S\) and \(T\) are finite and
      \(\operatorname{span}(S)=\operatorname{span}(T)\), then \(S\)
      and \(T\) have the same number of elements.
    \item \(\operatorname{span}(S\cap T)
      =\operatorname{span}(S)\cap\operatorname{span}(T)\).
    \item If \(S\subsetneq T\) and \(T\) is linearly independent, then
      \(\operatorname{span}(S)\subsetneq\operatorname{span}(T)\).
  \end{enumerate}
  \solutionlink{hw01:span-and-independence-assertions}
\end{exercise}

% !TEX root = ../../hw01.tex
\noindent\emph{Awaiting solution.}


\end{document}
